% Generated by paper/make_figures.py --fig tab-scale.
% Source: experiments/v4_protocol/v4_protocol_{real93v4,gen493v4,gen1085v4}.json
% Do not edit by hand; re-run the script so the numbers cannot drift from the source.
\begin{table}[t]
  \centering
  \small
  \caption{Data-scale ablation. The state-editing arm is the only arm whose data can be scaled, so it carries the experiment: the identical recipe, snapshot schedule and evaluation are used on the first 493 edited images and on all 1085. Every row is the arm-level joint detector+checkpoint estimate, so the selected detector is not the same in every row. The 93-image row is arm A on real photographs: it is the same protocol's real-data baseline, not a point on the synthetic scale, and no scale gain is claimed for it. The 493 $\rightarrow$ 1085 step is quantified in Table~\ref{tab:bootstrap}.}
  \label{tab:scale}
  \setlength{\tabcolsep}{4pt}
  \begin{tabular}{lccc}
  \toprule
  Training images & OOF mAP50--95 & OOF mAP50 & Deployment \\
  \midrule
  93 (real photographs) & 0.4320 & 0.6738 & yolo26s @ 280 \\
  493 (state-edited) & 0.5019 & 0.7367 & yolo26m @ 200 \\
  1085 (state-edited) & 0.6429 & 0.8370 & yolo26s @ 260 \\
  \bottomrule
  \end{tabular}
\end{table}
